% !TeX root = ../regression_logistique.tex
\chapter{Nuage de points et frontière de décision}
\label{ch:nuage}

\orient{%
écrire l'équation d'une frontière de décision à partir d'un vecteur normal et
d'un point de passage ; calculer un score et l'interpréter comme une distance
signée ; coder la variable cible}{%
produit scalaire, norme (\annexeref{ann:prerequis}) ;
chapitre~\ref{ch:standardisation}}{%
essentiel ; géométrie en dimension quelconque à l'\annexeref{ann:geometrie}}

\section{Séparation par une note}

\Herb sépare les maisons en deux groupes : \Gryf et \Slyth d'un côté,
\Huff et \Rav de l'autre. Sur cette seule note, un \Gryf et un
\Slyth sont indiscernables. \Runes les sépare autrement : \Gryf et
\Rav d'un côté, \Huff et \Slyth de l'autre. Là encore, deux maisons
restent confondues.

\begin{figure}[htbp]
\centering
\begin{graphique}{Deux histogrammes superposés de la note d'Herbology standardisée, de -3 à 3. Les trois autres maisons, hachurées, forment deux bosses. Les Gryffindor, en aplat bleu, se concentrent sous la bosse de gauche, partagée avec une partie des autres maisons.}
\begin{tikzpicture}
\begin{axis}[width=.78\textwidth,height=5.4cm,
  axis lines=left,axis line style={ink!35},
  grid=major,grid style={gridgray!35},
  tick label style={font=\small},label style={font=\small},
  xlabel={\Herb, en écarts-types},ylabel={proportion d'élèves},
  xmin=-3,xmax=3,ymin=0,ytick=\empty,
  legend style={at={(0.5,1.1)},anchor=south,legend columns=-1,draw=none,
    font=\small,/tikz/every even column/.append style={column sep=10pt}}]
  \addplot[hist={bins=34,density},pattern=north east lines,pattern color=warning!75,draw=warning!85]
    table[y=x] {data/autres.dat};
  \addlegendentry{les trois autres maisons}
  \addplot[hist={bins=34,density},fill=accent!55,draw=accent!80,opacity=.75]
    table[y=x] {data/gryffindor.dat};
  \addlegendentry{\Gryf}
\end{axis}
\end{tikzpicture}
\end{graphique}
\caption{Répartition selon la seule note d'\Herb ; trois autres maisons
hachurées, \Gryf en aplat. Les deux populations partagent le mode de gauche :
aucun seuil sur cet axe ne les sépare.}
\label{fig:une-note}
\end{figure}

Prises ensemble, les deux notes situent chaque élève dans un plan, où les \Gryf
occupent une région que les trois autres maisons laissent presque vide. La
séparation qu'aucune des deux notes ne produit seule résulte de leur
combinaison.

\begin{figure}[htbp]
\centering
\begin{graphique}{Nuage de points des deux notes standardisées, axes de -2,8 à 2,8, zéros en pointillé. Les Gryffindor, disques bleus, occupent le quadrant en haut à gauche : Herbology basse et Ancient Runes haute. Les trois autres maisons, croix orangées, occupent les trois autres quadrants, avec un léger recouvrement au centre.}
\begin{tikzpicture}
\begin{axis}[cloud,
  xlabel={\Herb, en écarts-types},ylabel={\Runes, en écarts-types},
  xmin=-2.8,xmax=2.8,ymin=-2.8,ymax=2.8]
  \draw[gridgray,dashed] (axis cs:-2.8,0)--(axis cs:2.8,0);
  \draw[gridgray,dashed] (axis cs:0,-2.8)--(axis cs:0,2.8);
  \addplot[only marks,mark=x,mark size=1.6pt,line width=.4pt,warning!85,opacity=.6]
    table[x=x,y=y] {data/autres.dat};
  \addlegendentry{les trois autres maisons}
  \addplot[only marks,mark=*,mark size=1.4pt,accent,opacity=.8]
    table[x=x,y=y] {data/gryffindor.dat};
  \addlegendentry{\Gryf}
\end{axis}
\end{tikzpicture}
\end{graphique}
\caption{Les mêmes élèves placés selon leurs deux notes. Le quadrant $x_1<0$,
$x_2>0$ contient $310$ élèves dont $304$ \Gryf ; les $23$ \Gryf restants se
trouvent ailleurs. Ce recouvrement borne ce qu'une droite peut atteindre.}
\label{fig:nuage}
\end{figure}
\FloatBarrier

\section{Équation de la frontière}

Les deux notes standardisées $(x_1,x_2)$ sont des nombres sans dimension,
exprimés dans la même unité --- l'écart-type de leur colonne. Le plan qu'elles
engendrent peut donc être muni d'un repère orthonormé, et c'est ce qui donne un
sens aux notions d'angle, de perpendiculaire et de distance employées plus bas.
Sur les notes brutes, où un axe se gradue en points et l'autre en centaines de
points, ces deux notions dépendraient du barème de chaque matière.

Une droite du plan se construit à partir de deux données : une direction
perpendiculaire et un point de passage. Notons $w=(w_1,w_2)$ un vecteur
perpendiculaire à la droite cherchée, et $a$ un point par lequel elle passe. Un
point $x$ appartient alors à la droite lorsque le vecteur $x-a$ est
perpendiculaire à $w$, c'est-à-dire lorsque leur produit scalaire s'annule :
\[
  w\cdot(x-a)=0,
  \qquad\text{où}\qquad
  w\cdot v=w_1v_1+w_2v_2 .
\]
Développer ce produit et poser $w_0=-w\cdot a$ donne l'équation sous sa forme
usuelle :
\[
  \ann{accent}{w_1x_1+w_2x_2}{\annl{direction}\\ \annl{perpendiculaire}}
  \;+\;\ann{warning}{w_0}{\annl{position}\\ \annl{de la droite}}\;=\;0
\]

\begin{figure}[htbp]
\centering
\begin{graphique}{Sur une grille, une droite oblique sépare le plan. Depuis un point a de la droite partent deux vecteurs : w, perpendiculaire à la droite et pointant vers le demi-plan supérieur, et x moins a, vers un point x. Un petit carré marque l'angle droit entre eux. Les deux demi-plans portent les mentions w·x + w0 positif et négatif.}
\begin{tikzpicture}[x=1.35cm,y=1.35cm,font=\footnotesize,
  >={Stealth[length=2.2mm]}]
\begin{scope}
\clip (-0.6,-0.5) rectangle (6.5,4.2);
\draw[gridgray!40,step=1] (-1,-1) grid (7,5);
\fill[accent!8] (-1,4.5) -- (7,0.5) -- (7,5) -- (-1,5) -- cycle;
\draw[ink,very thick] (-1,4.5) -- (7,0.5);
\end{scope}
\draw[ink!45] (-0.6,-0.5) rectangle (6.5,4.2);

\coordinate (a) at (2,3);
\coordinate (x) at (4.6,1.7);
\draw[warning,thick,->] (a) -- ++(0.894,1.789)
  node[anchor=south west,text=warning] {$w=(w_1,w_2)$};
\draw[accent,thick,->] (a) -- (x) node[midway,anchor=north west,text=accent,yshift=-1mm] {$x-a$};

% l'angle droit entre w et x-a
\draw[ink!55] ($(a)+(0.40,0.80)$) -- ($(a)+(0.40,0.80)+(0.775,-0.387)$)
  -- ($(a)+(0.775,-0.387)$);

\fill[ink] (a) circle (2pt) node[anchor=south east,inner sep=3pt] {$a$};
\fill[ink] (x) circle (2pt) node[anchor=north,inner sep=4pt] {$x$};
\node[accent!85,anchor=west] at (4.9,3.6) {$w\cdot x+w_0>0$};
\node[ink!70,anchor=west] at (0.1,0.35) {$w\cdot x+w_0<0$};
\node[ink,anchor=north west,fill=white,inner sep=1.5pt] at (5.15,1.42)
  {$w\cdot x+w_0=0$};
\end{tikzpicture}
\end{graphique}
\caption{La droite est l'ensemble des points $x$ tels que $x-a$ soit
perpendiculaire à $w$. Le vecteur $w$ pointe vers le demi-plan où l'expression
$w\cdot x+w_0$ est positive.}
\label{fig:construction}
\end{figure}
\FloatBarrier

Le couple $(w_1,w_2)$ est le \emph{vecteur normal} : perpendiculaire à la droite
par construction, il pointe vers le demi-plan où $w_0+w_1x_1+w_2x_2$ est positif,
puisque s'éloigner de la droite dans sa direction fait croître le produit
scalaire. Sa longueur, la \emph{norme} de $w$, intervient dans la distance à la
frontière et dans la raideur de la fonction logistique :
\[
  \|w\|=\sqrt{w_1^2+w_2^2}
\]

Le coefficient $w_0$ règle la position. À direction fixée, le faire varier
translate la droite parallèlement à elle-même, et la distance signée de l'origine
à la droite vaut $-w_0/\|w\|$. On l'appelle le \emph{biais} : c'est la valeur que
prend l'expression pour une observation dont les deux notes valent zéro.

Trois coefficients apparaissent alors que deux nombres suffisent à désigner une
droite du plan. La raison tient à la construction : multiplier $w$ et $w_0$ par
un même $\lambda\neq0$ multiplie tout le produit scalaire par $\lambda$ sans
changer les points où il s'annule. Seuls les rapports entre coefficients fixent
la droite.

\begin{approfondir}[Effet de chaque coefficient]
Tracés de la frontière lorsqu'un coefficient varie, les autres restant fixés :
\annexeref{ann:geometrie}, section~2.
\end{approfondir}

\section{Score}

Le membre de gauche de cette équation se calcule pour tout élève du plan, sur la
frontière ou ailleurs. On l'appelle le \emph{score} :
\[
  z=\ann{rulecolor}{w_0}{\annl{biais}\\ \annl{le même pour tous}}
   +\ann{accent}{w_1x_1}{\annl{apport de la note}\\ \annl{d'\Herb}}
   +\ann{warning}{w_2x_2}{\annl{apport de la note}\\ \annl{d'\Runes}}
\]
La frontière de décision est l'ensemble des points où ce nombre s'annule,
$\{x : z(x)=0\}$. Ailleurs, la construction de la section précédente en donne
directement le sens : puisque $w_0=-w\cdot a$, le score s'écrit
\[
  z=w\cdot x+w_0=w\cdot(x-a) .
\]
Or un produit scalaire vaut le produit des longueurs par le cosinus de l'angle,
et $\|x-a\|\cos\theta$ est exactement la distance de $x$ à la droite, comptée
perpendiculairement et affectée du signe du côté. En notant $d$ cette distance,
\[
  z=\|w\|\,d .
\]
Le signe du score donne donc le côté de la frontière, et sa valeur absolue la
distance, au facteur $\|w\|$ près. Deux observations de même score se tiennent à
la même distance de la frontière, du même côté. La conversion de ce score en
probabilité fait l'objet du chapitre~\ref{ch:probabilite}.

\begin{figure}[htbp]
\centering
\begin{graphique}{Un plan avec une droite frontière et deux demi-plans teintés. Trois points marqués portent leurs scores : +3,1 en haut à gauche, +0,2 près de la frontière, -2,7 en bas à droite. Un segment en pointillé joint le point de score +3,1 à son projeté orthogonal sur la frontière et mesure 2,66.}
\begin{tikzpicture}
% axis equal image : sans echelles identiques sur les deux axes, la
% perpendiculaire ne se lit pas comme telle.
\begin{axis}[width=.66\textwidth,axis equal image,
  axis lines=box,axis line style={ink!40},
  grid=major,grid style={gridgray!30},
  tick label style={font=\small},label style={font=\small},
  xlabel={$x_1$},ylabel={$x_2$},
  xmin=-3,xmax=3,ymin=-2.2,ymax=3,
  xtick={-3,-2,-1,0,1,2,3},ytick={-2,-1,0,1,2,3},
  domain=-3:3,samples=2,clip=true]
  \fill[accent!12] (axis cs:-3,-1.4) -- (axis cs:3,2.2)
    -- (axis cs:3,3) -- (axis cs:-3,3) -- cycle;
  \fill[warning!10] (axis cs:-3,-1.4) -- (axis cs:3,2.2)
    -- (axis cs:3,-2.2) -- (axis cs:-3,-2.2) -- cycle;
  \addplot[ink,very thick]{0.4+0.6*x};
  \addplot[ink!55,thin,dashed] coordinates {(-1.5,2.6) (-0.132,0.321)};
  \addplot[only marks,mark=*,ink,mark size=2.2pt]
    coordinates {(-1.5,2.6) (0.5,0.9) (1.5,-1.4)};
  \node[font=\small,text=ink,anchor=south west] at (axis cs:-1.45,2.68)
    {$z=+3{,}1$};
  \node[font=\small,text=ink,anchor=south] at (axis cs:0.5,1.02) {$z=+0{,}2$};
  \node[font=\small,text=ink,anchor=west] at (axis cs:1.62,-1.4) {$z=-2{,}7$};
  \node[font=\small,text=ink,fill=white,inner sep=1.5pt,anchor=north west]
    at (axis cs:2.2,1.72) {$z=0$};
  \node[font=\small,text=accent,anchor=north west] at (axis cs:-2.9,2.9)
    {$z>0$};
  \node[font=\small,text=warning,anchor=south east] at (axis cs:2.9,-2.1)
    {$z<0$};
  \node[font=\scriptsize,text=ink!75,anchor=west]
    at (axis cs:-0.75,1.50) {$d=2{,}66$};
\end{axis}
\end{tikzpicture}
\end{graphique}
\caption{Frontière pour $w=(-0{,}4\,;-0{,}6\,;1)$. Les distances orthogonales
des points de scores $+3{,}1$ et $+0{,}2$ valent respectivement $2{,}66$ et
$0{,}17$.}
\label{fig:score}
\end{figure}

\section{Variable cible}

La frontière de décision sépare le plan en exactement deux régions. Elle peut
donc porter une question à deux réponses, et une seule, alors que le fichier en
distingue quatre.

La construction traite donc une maison à la fois : \Gryf contre les trois autres
réunies. Les trois questions restantes se posent à l'identique et se résolvent
séparément (ch.~\ref{ch:classifieur}).

La réponse à cette question s'appelle la \emph{variable cible}. Sa colonne se
déduit ligne à ligne de la colonne \texttt{Hogwarts House} du fichier, et c'est
la quantité que le modèle devra produire à partir des seules notes :
\[
  y_i=\begin{cases}
    1 & \text{si l'observation } i \text{ appartient à \Gryf},\\
    0 & \text{sinon,}
  \end{cases}
\]
soit $327$ valeurs égales à $1$ et $1\,273$ égales à $0$.

Ce codage rend la moyenne de la colonne $y$ égale à la proportion de \Gryf,
$327/1\,600=0{,}2044$, et la rend comparable ligne à ligne à la colonne des
probabilités produite au chapitre suivant : leur différence mesure l'écart entre
annoncé et observé.

\section{Propriétés requises de la transformation}

Le signe de $z$ fournit déjà une règle de décision : \Gryf si $z>0$, autre maison
sinon. Deux exigences conduisent à lui préférer une quantité qui varie de façon
continue.

La première est de rendre compte de la fiabilité. Deux observations situées du
même côté reçoivent la même réponse, que l'une se tienne loin de la frontière et
l'autre tout contre, là où une note manquante ou une erreur de saisie suffit à la
faire changer de classe. La valeur du score porte cette différence, absente de
son signe : les deux points marqués sur la figure~\ref{fig:score} ont pour scores
$+3{,}1$ et $+0{,}2$.

La seconde est de rendre la frontière corrigeable. Déterminer sa position par le
calcul suppose de mesurer de combien elle est mal placée, de la déplacer
d'autant, et de recommencer. Une réponse binaire garde la même valeur sur tout un
intervalle de positions, et le calcul y reçoit partout la même indication. Il
faut une quantité qui réponde à un déplacement, si petit soit-il.

\begin{resultat}[Transformation cherchée]
Transformer le score $z$, qui parcourt $\R$, en une probabilité comprise entre
$0$ et $1$, par une fonction continue --- un petit déplacement de la droite
produit un petit changement de probabilité --- et croissante.

Une infinité de fonctions remplissent ces deux conditions : la fonction de
répartition de la loi normale, qui fonde le modèle \emph{probit}, et la fonction
logistique, retenue ici. Ce choix est une décision de modélisation
(\annexeref{ann:logit}).
\end{resultat}

\enlargethispage{2\baselineskip}
Le modèle complet du projet retient dix matières : une observation est alors un
point de $\R^n$ et la frontière un \emph{hyperplan affine} de dimension $n-1$.
Le vecteur normal, la distance $|z|/\|w\|$ et le signe de $z$ gardent leur rôle ;
l'équation et le décompte des coefficients en dimensions $1$, $2$ et $3$ sont
détaillés dans l'\annexeref{ann:geometrie}. Le code de la partie~II travaille
sur la liste \texttt{COURSES}, dont la longueur fixe le nombre de matières.

\begin{exercice}[Vecteur normal et distance]
Pour $w=(-4{,}745\,;-3{,}454\,;+3{,}469)$, calculer $\|w\|$, la distance signée
de l'origine à la frontière, et le score de l'observation $x=(0\,;0)$. Vérifier
que le signe obtenu place l'élève moyen du côté attendu au vu de la
figure~\ref{fig:nuage}.
\end{exercice}
